\section{Finite Graph Tensor Networks}\label{sec:peps_foundations}

We establish the tensor-network setting on finite graphs: the local tensors,
the state they generate, the gauge freedom on the bonds, and the vertex
injectivity condition together with the local inverse it provides.

\subsection{PEPS tensors and the states they generate}

\begin{definition}[Graph edge for PEPS]
    \label{def:peps_edge}
    \lean{TNLean.PEPS.Edge}
    \leanok
    For a finite simple graph $G$, an \emph{edge} is an ordered pair
    $(u,v)$ with $u < v$ and $u$ adjacent to $v$.
    The ordering serves to avoid double counting.
\end{definition}

\begin{definition}[Incident edges at a vertex]
    \label{def:peps_incident_edge}
    \lean{TNLean.PEPS.IncidentEdge}
    \leanok
    \uses{def:peps_edge}
    For a vertex $v$, an \emph{incident edge} is an edge of $G$ having $v$
    as one of its endpoints.
\end{definition}

\begin{definition}[PEPS tensor on a graph]
    \label{def:peps_tensor}
    \lean{TNLean.PEPS.Tensor}
    \leanok
    \uses{def:peps_edge, def:peps_incident_edge}
    A PEPS tensor $A$ on $G=(V,E)$ with physical dimension $d$ consists of a bond
    dimension $D_e$ for each $e\in E$ and a local tensor
    $A_v:\prod_{e\ni v}\{0,\ldots,D_e-1\}\times\{0,\ldots,d-1\}\to\C$ at each
    $v\in V$, with one physical leg and one virtual leg per incident edge:
    % Figure peps_tensor.
% Owner: chapter/ch24_peps_ft_foundations.tex (def:peps_tensor).
% Used by: gapped-2d-notes PP Sec. 0 (PEPS definition).
% Formula: A_v in (tensor_{e incident to v} C^{D_e})^* tensor C^d.
% Ink-to-index: the three drawn thin legs represent the first, a typical,
% and the last member of the arbitrary incident-edge family; intervening
% incident edges are suppressed, and the north leg is sigma(v).
% Contraction and boundary signature: no index is contracted; the exact
% tensor has (V,P)=(deg(v),1), independently of the three-leg schematic.
% Source verdict: arbitrary-graph local tensor of
% Papers/1804.04964/paper_normal.tex:961--979; suppressing intervening
% edges avoids a square-lattice degree-four restriction.
% Related: the injectivity-map display of the same chapter
% (def:IsVertexInjective, ker A_v = 0) is not reused: injectivity is not
% a hypothesis of the PEPS note's Theorem 1.
\begin{center}
    \tenkzkernel
    \begin{tenkz}[rows={wire}, cols=1, bonds=none]
        \tn[label pos=w,
            ports={210:virtual:$e_1$, 270:virtual:$e_k$,
                    330:virtual:$e_{\deg(v)}$,
                    90:physical:$\sigma(v)$}]{A_v}
    \end{tenkz}
\end{center}

\end{definition}

\begin{definition}[Virtual configuration]
    \label{def:peps_virtual_config}
    \lean{TNLean.PEPS.VirtualConfig}
    \leanok
    \uses{def:peps_tensor, def:peps_edge}
    A \emph{virtual configuration} is an assignment
    $\eta\in\prod_{e\in E}\{0,\ldots,D_e-1\}$.
\end{definition}

\begin{definition}[Splitting and reindexing virtual configurations]
    \label{def:peps_virtual_config_split_reindex}
    \lean{TNLean.PEPS.virtualConfigSplit}
    \lean{TNLean.PEPS.IsIncidentEdge}
    \lean{TNLean.PEPS.vertexConfigSplitAt}
    \lean{TNLean.PEPS.virtualConfigCongr}
    \leanok
    \uses{def:peps_virtual_config}
    An edge predicate $p$ splits a virtual configuration $\eta$ into
    $(\eta|_p,\eta|_{\neg p})$. A bond-dimension equality $h:D=D'$ induces a
    canonical reindexing $\iota_h$ of virtual configurations. For $v\in V$,
    selecting incident edges splits $\eta$ into $(\eta|_v,\eta|_{\not\ni v})$.
\end{definition}

\begin{theorem}[Evaluation laws for splitting and reindexing]
    \label{thm:peps_virtual_config_split_reindex_apply}
    \lean{TNLean.PEPS.virtualConfigSplit_fst_apply}
    \lean{TNLean.PEPS.virtualConfigSplit_snd_apply}
    \lean{TNLean.PEPS.isIncidentEdge_of_incident}
    \lean{TNLean.PEPS.vertexConfigSplitAt_fst}
    \lean{TNLean.PEPS.vertexConfigSplitAt_symm_apply_incident}
    \lean{TNLean.PEPS.virtualConfigCongr_apply}
    \lean{TNLean.PEPS.virtualConfigSplit_mergeVirtualConfig}
    \lean{TNLean.PEPS.virtualConfigCongr_mergeVirtualConfig}
    \leanok
    \uses{def:peps_virtual_config_split_reindex, def:peps_virtual_config_merge}
    On their respective domains, $(\eta|_p)(e)=\eta(e)$ and
    $(\eta|_{\neg p})(e)=\eta(e)$. With $m_p(\eta,\xi)$ the piecewise merge of
    Definition~\ref{def:peps_virtual_config_merge},
    $(\iota_h\eta)(e)=h_e(\eta(e))$,
    $\iota_h(m_p(\eta,\xi))=m_p(\iota_h\eta,\iota_h\xi)$, and
    $(m_p(\eta,\xi)|_p,m_p(\eta,\xi)|_{\neg p})=(\eta|_p,\xi|_{\neg p})$.
\end{theorem}
\begin{proof}\leanok
    Evaluate at $e\in E$ and split by whether $p(e)$ holds.
\end{proof}

\begin{definition}[Piecewise and regionwise merge of virtual configurations]
    \label{def:peps_virtual_config_merge}
    \lean{TNLean.PEPS.mergeVirtualConfig}
    \lean{TNLean.PEPS.regionMerge}
    \leanok
    \uses{def:peps_virtual_config}
    For virtual configurations $\eta,\xi$ and an edge predicate $p$, the
    \emph{piecewise merge} $m_p(\eta,\xi)$ is defined by
    $m_p(\eta,\xi)(e)=\eta(e)$ if $p(e)$ and $\xi(e)$ otherwise. For a region
    $R\subseteq V$, the \emph{regionwise merge} $m_R(\eta,\xi)$ takes $p(e)$ to
    hold iff $e$ is incident to $R$.
\end{definition}

\begin{theorem}[Readback laws for merged configurations]
    \label{thm:peps_virtual_config_merge_readback}
    \lean{TNLean.PEPS.mergeVirtualConfig_of_pos}
    \lean{TNLean.PEPS.mergeVirtualConfig_of_neg}
    \lean{TNLean.PEPS.mergeVirtualConfig_self}
    \lean{TNLean.PEPS.regionMerge_of_incident}
    \lean{TNLean.PEPS.regionMerge_of_not_incident}
    \leanok
    \uses{def:peps_virtual_config_merge}
    Under Definition~\ref{def:peps_virtual_config_merge},
    $m_p(\eta,\xi)(e)=\eta(e)$ if $p(e)$, $m_p(\eta,\xi)(e)=\xi(e)$ if
    $\neg p(e)$, and $m_p(\eta,\eta)=\eta$. In particular,
    $m_R(\eta,\xi)(e)=\eta(e)$ for $e$ incident to $R$ and $\xi(e)$ otherwise.
\end{theorem}
\begin{proof}\leanok
    Evaluate $m_p(\eta,\xi)(e)$ by case distinction on $p(e)$.
\end{proof}

\begin{theorem}[Vertex products through a regionwise merge]
    \label{thm:peps_region_merge_products}
    \lean{TNLean.PEPS.regionProd_eq_merge}
    \lean{TNLean.PEPS.regionProd_p2_eq_merge_of_incident_agree}
    \lean{TNLean.PEPS.complementProd_eq_merge}
    \leanok
    \uses{def:peps_virtual_config_merge}
    Put $\Pi_R(\eta,\sigma)=\prod_{v\in R}A_v(\eta|_v,\sigma_v)$. Then
    $\Pi_R(m_R(\eta,\xi),\sigma)=\Pi_R(\eta,\sigma)$, and if $\eta(e)=\xi(e)$
    for every $e$ incident to both $R$ and $S$, then
    $\Pi_S(m_R(\eta,\xi),\sigma)=\Pi_S(\xi,\sigma)$. In particular, if
    $\eta=\xi$ on $\partial R$, then
    $\Pi_{V\setminus R}(m_R(\eta,\xi),\sigma)=\Pi_{V\setminus R}(\xi,\sigma)$.
\end{theorem}
\begin{proof}\leanok
    By Theorem~\ref{thm:peps_virtual_config_merge_readback},
    $m_R(\eta,\xi)|_v=\eta|_v$ for $v\in R$, while for $v\in S$ any edge $e\ni v$
    also incident to $R$ satisfies $\eta(e)=\xi(e)$; taking $S=V\setminus R$
    specializes this to $\partial R$.
\end{proof}

\begin{theorem}[Constant-cardinality fiber sum]
    \label{thm:peps_config_fiber_sum}
    \lean{TNLean.PEPS.sum_comp_of_config_fiber_card}
    \leanok
    \uses{def:peps_virtual_config}
    If $m:X\to\prod_e\{0,\ldots,D_e-1\}$ has fibers $|m^{-1}(\eta)|=n$ for all
    $\eta$, then for every function $g$ into an additive commutative monoid,
    $\sum_{x\in X} g(m(x))=n\mathbin{\boldsymbol{\cdot}}\sum_{\eta}g(\eta)$.
\end{theorem}
\begin{proof}\leanok
    Partitioning $X$ into the fibers $m^{-1}(\eta)$ repeats each summand
    $g(\eta)$ exactly $|m^{-1}(\eta)|=n$ times.
\end{proof}

\begin{definition}[PEPS state coefficient]
    \label{def:peps_stateCoeff}
    \lean{TNLean.PEPS.stateCoeff}
    \leanok
    \uses{def:peps_tensor, def:peps_virtual_config}
    For a physical configuration $\sigma:V\to\{0,\ldots,d-1\}$, the PEPS state
    coefficient is
    $\Coeff_A(\sigma)=\sum_{\eta}\prod_{v\in V}A_v(\eta|_v,\sigma_v)$.
\end{definition}

\begin{definition}[Equality of PEPS states]
    \label{def:peps_same_state}
    \lean{TNLean.PEPS.SameState}
    \leanok
    \uses{def:peps_stateCoeff}
    Two PEPS tensors $A,B$ represent the same state if
    $\Coeff_A(\sigma)=\Coeff_B(\sigma)$ for all $\sigma$.
\end{definition}

\subsection{Gauge transformations on the bonds}

An invertible matrix $X_e\in\GL(D_e)$ on an edge acts on its two endpoint
tensors by opposite oriented insertions; absorbing $X_e$ on every edge gives a
gauge transformation preserving $\Coeff_A$.

\begin{definition}[Edge gauge at a vertex]%
    \label{def:edgeGaugeAt}
    \lean{TNLean.PEPS.edgeGaugeAt}
    \leanok
    For an edge $e=(u,w)$ with $u<w$ and $X_e\in\GL(D_e)$, the endpoint gauge
    matrix at $v\in\{u,w\}$ is
    % Formula: the two endpoint actions of the ordered edge e=(u,w) are X_e and
    % (X_e^{-1})^top.  No standalone source figure states this convention.
    % Ink-to-index/type verdict: the old scalar picture suppressed the physical and
    % external virtual legs; the actual two-vertex object would have boundary
    % (V,P)=(deg(u)+deg(w)-2,2).  The exact convention is therefore kept as algebra.
    \begin{align}
        X_{v,e}
         & =
        \begin{cases}
            X_e,             & v=u, \\
            (X_e^{-1})^\top, & v=w.
        \end{cases}
        \notag
    \end{align}
\end{definition}

\begin{definition}[Gauge vertex action]%
    \label{def:gaugeVertex}
    \lean{TNLean.PEPS.gaugeVertex}
    \leanok
    \uses{def:edgeGaugeAt}
    Writing $M_{ie}=X_{v,e}$ from Definition~\ref{def:edgeGaugeAt} for each
    incident edge $ie=(v,e)$, the gauged vertex tensor $(A^X)_v$ is
    % Formula: the displayed component equation is the gauge action at v.
    % Ink-to-index/type verdict: Papers/1804.04964/paper_normal.tex:1212--1314
    % draws only particular vertices.  The general object has boundary
    % (V,P)=(deg(v),1), whereas the old picture hard-coded degree four; retain the
    % valence-independent equation rather than asserting a universal star picture.
    \begin{align}
        (A^X)_v(\eta,\sigma)
         & =
        \sum_{\eta'}
        \left(\prod_{ie\ni v}
        M_{ie}(\eta(ie),\eta'(ie))\right)A_v(\eta',\sigma).
        \notag
    \end{align}
\end{definition}

\begin{definition}[Apply gauge]%
    \label{def:applyGauge}
    \lean{TNLean.PEPS.applyGauge}
    \leanok
    \uses{def:gaugeVertex}
    Applying $X\in\prod_e\GL(D_e)$ to $A$ yields the PEPS tensor $A^X$ with
    vertex tensors $(A^X)_v$ from Definition~\ref{def:gaugeVertex}.
\end{definition}

\begin{definition}[Absorb edge gauges]%
    \label{def:peps_absorbEdgeGauges}
    \lean{TNLean.PEPS.absorbEdgeGauges}
    \leanok
    \uses{def:applyGauge}
    For a tensor family $B$ and edge gauges $X=(X_e)_e$, the absorbed tensor
    family is $\widetilde B=B^X$.
\end{definition}

\begin{theorem}[Absorbed gauges keep the bond spaces]%
    \label{thm:peps_absorbEdgeGauges_bondDim}
    \lean{TNLean.PEPS.absorbEdgeGauges_bondDim}
    \leanok
    \uses{def:peps_absorbEdgeGauges}
    For every edge $e$, $D_{\widetilde B}(e)=D_B(e)$.
\end{theorem}
\begin{proof}\leanok
    By Definition~\ref{def:peps_absorbEdgeGauges},
    $D_{\widetilde B}(e)=D_{B^X}(e)=D_B(e)$.
\end{proof}

\begin{theorem}[Local formula for absorbed gauges]%
    \label{thm:peps_absorbEdgeGauges_component}
    \lean{TNLean.PEPS.absorbEdgeGauges_component}
    \leanok
    \uses{def:peps_absorbEdgeGauges, def:gaugeVertex}
    With $M_{ie}=X_{v,e}$ on each $ie\ni v$,
    $\widetilde B_v(\eta,\sigma)
        =\sum_{\eta'}(\prod_{ie\ni v}M_{ie}(\eta(ie),\eta'(ie)))
        B_v(\eta',\sigma)$.
\end{theorem}
\begin{proof}\leanok
    Unfold $\widetilde B_v=(B^X)_v$ using Definition~\ref{def:gaugeVertex}.
\end{proof}

\begin{definition}[Gauge equivalence for PEPS]%
    \label{def:GaugeEquivPEPS}
    \lean{TNLean.PEPS.GaugeEquiv}
    \leanok
    \uses{def:gaugeVertex}
    Two PEPS tensors $A,B$ are gauge-equivalent ($A\sim B$) if $D_A=D_B$ and
    $B=A^X$ for some $X\in\prod_e\GL(D_e)$.
\end{definition}

\begin{theorem}[Gauge invariance of PEPS state]%
    \label{thm:applyGauge_stateCoeff}
    \lean{TNLean.PEPS.applyGauge_stateCoeff}
    \leanok
    \uses{def:applyGauge, def:peps_stateCoeff}
    For every $\sigma$, $\Coeff_{A^X}(\sigma)=\Coeff_A(\sigma)$, since on each
    edge $e=(u,w)$ the endpoint insertions satisfy
    % Formula: contraction of X_e with the transposed inverse on the other endpoint
    % gives the Kronecker delta.  Ink-to-index/type verdict: i and k remain free, so
    % this identity kernel has (V,P)=(2,0).  It has no separate source figure; the
    % old two-star drawing introduced unrelated external legs.
    \begin{align}
        \sum_j X_e(j,i)((X_e^{-1})^\top)(j,k)
         & =\delta_{ik}.
        \notag
    \end{align}
\end{theorem}

\begin{theorem}[Gauge equivalence implies same state]%
    \label{thm:GaugeEquiv_sameState}
    \lean{TNLean.PEPS.GaugeEquiv.sameState}
    \leanok
    \uses{def:GaugeEquivPEPS, def:peps_same_state, thm:applyGauge_stateCoeff}
    If $A \sim B$, then $\Coeff_A(\sigma) = \Coeff_B(\sigma)$ for all $\sigma$.
\end{theorem}

\subsection{Vertex injectivity and the local inverse}

When the local tensor map $\Phi_v:\eta_v\mapsto A_v(\eta_v,\cdot)\in\C^d$ is
injective, a left inverse $\Psi_v$ converts virtual bond insertions at $v$ into
physical operators on $\C^d$.

\begin{definition}[Vertex injectivity]%
    \label{def:IsVertexInjective}
    \lean{TNLean.PEPS.IsVertexInjective}
    \leanok
    A PEPS tensor $A$ is \emph{vertex-injective} if at every $v\in V$ the family
    $(A_v(\eta,\cdot)\in\C^d)_{\eta}$ indexed by local virtual configurations
    $\eta\in\mathcal V_v=\prod_{e\ni v}\{0,\ldots,D_e-1\}$ is linearly
    independent in $\C^d$, i.e., its linear extension has $\ker A_v=\{0\}$ as in
    \cite[Section~3]{Molnar2018NormalPEPS}.
    % Formula: A_v : tensor_{e incident to v} C^{D_e} -> C^d and ker(A_v)=0.
    % Ink-to-index: the west wire bundles the virtual indices on every edge
    % incident to v; the east wire is the physical vector A_v(eta,cdot).
    % Contraction and boundary signature: no index is summed in this local map;
    % its bundled signature is (V,P)=(1,1), for arbitrary deg(v).
    % Source verdict: exact injectivity map of
    % Papers/1804.04964/paper_normal.tex:961--981; no square-lattice valence is
    % assumed.
    \begin{center}
        \tenkzkernel
        % The domain and the codomain name the two open ends of the map, so
        % each is the tip label of the port whose type it belongs to.
        \begin{tenkz}[rows={wire}, cols=1, bonds=none]
            \tn[skin=box,
                ports={180:virtual:$\bigotimes_{e\ni v}\C^{D_e}$,
                        0:physical:$\C^d$}]{A_v}
        \end{tenkz}
        \qquad $\ker A_v=\{0\}$.
    \end{center}
\end{definition}

\begin{definition}[Local tensor map]%
    \label{def:peps_localTensorMap}
    \lean{TNLean.PEPS.localTensorMap}
    \leanok
    \uses{def:peps_tensor, def:peps_incident_edge}
    At $v\in V$, the local tensor map $\Phi_v:\C^{\mathcal V_v}\to\C^d$ is
    $\Phi_v(c)(\tau)=\sum_{\eta\in\mathcal V_v}c(\eta)A_v(\eta,\tau)$.
\end{definition}

\begin{lemma}[Trivial kernel in explicit coordinates]%
    \label{lem:peps_localCoeff_eq_zero_of_contract_zero}
    \lean{TNLean.PEPS.IsVertexInjective.localCoeff_eq_zero_of_contract_zero}
    \leanok
    \uses{def:IsVertexInjective, def:peps_localTensorMap}
    If $A$ is vertex-injective and $\Phi_v(R)(\tau)=\sum_{\eta}R(\eta)A_v(\eta,\tau)=0$
    for all $\tau$, then $R=0$.
\end{lemma}
\begin{proof}\leanok
    Since $\Phi_v(R)=0$ and $\ker\Phi_v=\{0\}$ by
    Definition~\ref{def:IsVertexInjective}, $R=0$.
\end{proof}

\begin{definition}[Local left inverse]%
    \label{def:peps_localLeftInverse}
    \lean{TNLean.PEPS.localLeftInverseAt}
    \leanok
    \uses{def:peps_localTensorMap}
    When $\Phi_v$ is injective, choose a linear map
    $\Psi_v:\C^d\to\C^{\mathcal V_v}$ with $\Psi_v\circ\Phi_v=\Id$.
\end{definition}

\begin{definition}[Physical realization of a local virtual operator]%
    \label{def:peps_physRealizeLocalOp}
    \lean{TNLean.PEPS.physRealizeLocalOpAt}
    \leanok
    \uses{def:peps_localTensorMap, def:peps_localLeftInverse}
    For $T\in\End(\C^{\mathcal V_v})$, define the physical realization
    $O_T=\Phi_v\circ T\circ\Psi_v\in\End(\C^d)$.
\end{definition}

\begin{theorem}[Local virtual operations on the image]%
    \label{thm:peps_physRealizeLocalOp_spec}
    \lean{TNLean.PEPS.physRealizeLocalOpAt_spec}
    \leanok
    \uses{def:peps_physRealizeLocalOp}
    For every $c\in\C^{\mathcal V_v}$, $O_T(\Phi_v(c))=\Phi_v(Tc)$.
\end{theorem}
\begin{proof}\leanok
    By Definitions~\ref{def:peps_localLeftInverse}
    and~\ref{def:peps_physRealizeLocalOp},
    $O_T\Phi_v=\Phi_vT\Psi_v\Phi_v=\Phi_vT$.
\end{proof}

\begin{theorem}[Injectivity of local physical realization]%
    \label{thm:peps_physRealizeLocalOp_injective}
    \lean{TNLean.PEPS.physRealizeLocalOpAt_injective}
    \leanok
    \uses{def:peps_physRealizeLocalOp, thm:peps_physRealizeLocalOp_spec, def:peps_localTensorMap}
    If $\Phi_v$ is injective and $O_T=O_{T'}$, then $T=T'$.
\end{theorem}
\begin{proof}\leanok
    By Theorem~\ref{thm:peps_physRealizeLocalOp_spec},
    $\Phi_v(Tc)=O_T(\Phi_v(c))=O_{T'}(\Phi_v(c))=\Phi_v(T'c)$, so injectivity
    of $\Phi_v$ gives $T=T'$.
\end{proof}

\begin{definition}[Local projector]%
    \label{def:peps_localProjector}
    \lean{TNLean.PEPS.localProjectorAt}
    \leanok
    \uses{def:peps_physRealizeLocalOp}
    The local projector is $P_v=O_{\Id}=\Phi_v\circ\Psi_v$.
\end{definition}

\begin{definition}[Virtual pullback of a local physical operator]%
    \label{def:peps_localVirtualOpOfPhysicalOp}
    \lean{TNLean.PEPS.localVirtualOpOfPhysicalOpAt}
    \leanok
    \uses{def:peps_localTensorMap, def:peps_localLeftInverse}
    In the physical-to-virtual step of \cite[Section~3]{Molnar2018NormalPEPS},
    for $O\in\End(\C^d)$ at $v$, define its virtual pullback
    $T_O=\Psi_v\circ O\circ\Phi_v\in\End(\C^{\mathcal V_v})$.
\end{definition}

\begin{theorem}[Virtual pullback realizes the projected physical action]%
    \label{thm:peps_localVirtualOpOfPhysicalOp_spec}
    \lean{TNLean.PEPS.localVirtualOpOfPhysicalOpAt_spec}
    \leanok
    \uses{def:peps_localVirtualOpOfPhysicalOp, def:peps_physRealizeLocalOp, def:peps_localProjector}
    For every $c\in\C^{\mathcal V_v}$,
    $\Phi_v(T_Oc)=P_v(O(\Phi_v(c)))$.
\end{theorem}
\begin{proof}\leanok
    By Definitions~\ref{def:peps_localProjector}
    and~\ref{def:peps_localVirtualOpOfPhysicalOp},
    $\Phi_vT_O=\Phi_v\Psi_vO\Phi_v=P_vO\Phi_v$.
\end{proof}

\begin{theorem}[Virtual pullback for image-preserving physical actions]%
    \label{thm:peps_localVirtualOpOfPhysicalOp_realizes_of_projector}
    \lean{TNLean.PEPS.localVirtualOpOfPhysicalOpAt_realizes_of_projector}
    \leanok
    \uses{thm:peps_localVirtualOpOfPhysicalOp_spec}
    If $O(\operatorname{im}\Phi_v)\subseteq\operatorname{im}\Phi_v$, then
    $\Phi_v(T_Oc)=O(\Phi_v(c))$ for all $c\in\C^{\mathcal V_v}$.
\end{theorem}
\begin{proof}\leanok
    Since $P_v|_{\operatorname{im}\Phi_v}=\Id$,
    Theorem~\ref{thm:peps_localVirtualOpOfPhysicalOp_spec} gives
    $\Phi_v(T_Oc)=P_v(O(\Phi_v(c)))=O(\Phi_v(c))$.
\end{proof}

\begin{theorem}[Recovery of a virtual operation from a physical realization]%
    \label{thm:peps_localVirtualOpOfPhysicalOp_eq_of_realizes}
    \lean{TNLean.PEPS.localVirtualOpOfPhysicalOpAt_eq_of_realizes}
    \leanok
    \uses{def:peps_localTensorMap, thm:peps_localVirtualOpOfPhysicalOp_spec}
    Suppose $\Phi_v$ is injective. If $O(\Phi_v(c))=\Phi_v(Tc)$ for all
    $c\in\C^{\mathcal V_v}$, then $T_O=T$.
\end{theorem}
\begin{proof}\leanok
    By Theorem~\ref{thm:peps_localVirtualOpOfPhysicalOp_spec},
    $\Phi_v(T_Oc)=P_v(\Phi_v(Tc))=\Phi_v(Tc)$, and injectivity of $\Phi_v$
    gives $T_O=T$.
\end{proof}

\begin{theorem}[Extensionality of physical-to-virtual pullback]%
    \label{thm:peps_localVirtualOpOfPhysicalOp_eq_of_projected_action_eq}
    \lean{TNLean.PEPS.localVirtualOpOfPhysicalOpAt_eq_iff_projected_action_eq}
    \leanok
    \uses{def:peps_localTensorMap, thm:peps_localVirtualOpOfPhysicalOp_spec}
    When $\Phi_v$ is injective, two physical operations $O,O'\in\End(\C^d)$
    satisfy
    \begin{align}
        T_O=T_{O'}
        \quad\Longleftrightarrow\quad
        \forall c\in\C^{\mathcal V_v},\ P_vO\Phi_v(c)=P_vO'\Phi_v(c).
        \notag
    \end{align}
\end{theorem}
\begin{proof}\leanok
    By Theorem~\ref{thm:peps_localVirtualOpOfPhysicalOp_spec},
    $P_vO\Phi_v(c)=\Phi_v(T_Oc)$ and $P_vO'\Phi_v(c)=\Phi_v(T_{O'}c)$;
    injectivity of $\Phi_v$ gives the equivalence.
\end{proof}

\begin{theorem}[Pullback of the canonical physical realization]%
    \label{thm:peps_localVirtualOpOfPhysicalOp_physRealizeLocalOp}
    \lean{TNLean.PEPS.localVirtualOpOfPhysicalOpAt_physRealizeLocalOpAt}
    \leanok
    \uses{def:peps_physRealizeLocalOp, thm:peps_physRealizeLocalOp_spec, thm:peps_localVirtualOpOfPhysicalOp_eq_of_realizes}
    When $\Phi_v$ is injective, $T_{O_T}=T$.
\end{theorem}
\begin{proof}\leanok
    Combine Theorem~\ref{thm:peps_physRealizeLocalOp_spec} with
    Theorem~\ref{thm:peps_localVirtualOpOfPhysicalOp_eq_of_realizes} for
    $O=O_T$.
\end{proof}

\begin{theorem}[Physical realization of a pulled-back physical operator]%
    \label{thm:peps_physRealizeLocalOp_localVirtualOpOfPhysicalOp}
    \lean{TNLean.PEPS.physRealizeLocalOpAt_localVirtualOpOfPhysicalOpAt}
    \leanok
    \uses{def:peps_physRealizeLocalOp, def:peps_localVirtualOpOfPhysicalOp, def:peps_localProjector}
    With $P_v=\Phi_v\circ\Psi_v$ and $T_O=\Psi_v\circ O\circ\Phi_v$,
    $O_{T_O}=\Phi_v\circ T_O\circ\Psi_v=P_v\circ O\circ P_v$.
\end{theorem}
\begin{proof}\leanok
    Unfold Definitions~\ref{def:peps_physRealizeLocalOp},
    \ref{def:peps_localProjector}, and~\ref{def:peps_localVirtualOpOfPhysicalOp}.
\end{proof}

\begin{theorem}[Projected local recovery]%
    \label{thm:peps_localVirtualOpOfPhysicalOp_eq_of_projected_realization_eq}
    \lean{TNLean.PEPS.localVirtualOpOfPhysicalOpAt_eq_of_projected_realization_eq}
    \leanok
    \uses{thm:peps_physRealizeLocalOp_injective, thm:peps_physRealizeLocalOp_localVirtualOpOfPhysicalOp}
    When $\Phi_v$ is injective, $P_vOP_v=O_T$ implies $T_O=T$.
\end{theorem}
\begin{proof}\leanok
    By Theorem~\ref{thm:peps_physRealizeLocalOp_localVirtualOpOfPhysicalOp},
    $O_{T_O}=P_vOP_v=O_T$, so Theorem~\ref{thm:peps_physRealizeLocalOp_injective}
    yields $T_O=T$.
\end{proof}

\begin{theorem}[Projected local recovery equivalence]%
    \label{thm:peps_localVirtualOpOfPhysicalOp_eq_iff_projected_realization_eq}
    \lean{TNLean.PEPS.localVirtualOpOfPhysicalOpAt_eq_iff_projected_realization_eq}
    \leanok
    \uses{thm:peps_physRealizeLocalOp_localVirtualOpOfPhysicalOp, thm:peps_localVirtualOpOfPhysicalOp_eq_of_projected_realization_eq}
    When $\Phi_v$ is injective, $P_vOP_v=O_T\iff T_O=T$.
\end{theorem}
\begin{proof}\leanok
    Combine Theorems~\ref{thm:peps_physRealizeLocalOp_localVirtualOpOfPhysicalOp}
    and~\ref{thm:peps_localVirtualOpOfPhysicalOp_eq_of_projected_realization_eq}.
\end{proof}

\begin{definition}[Matrix action on one incident virtual edge]%
    \label{def:peps_localIncidentMatrixOp}
    \lean{TNLean.PEPS.localIncidentMatrixOp}
    \leanok
    \uses{def:peps_incident_edge, def:peps_localVirtualConfigSplitAt}
    Under the split $\eta\leftrightarrow(i,\eta_{\widehat e})$ at an incident
    edge $e\ni v$ from Definition~\ref{def:peps_localVirtualConfigSplitAt}, a
    matrix $M\in\MN{D_e}$ defines $T_{v,e}(M)\in\End(\C^{\mathcal V_v})$ by
    $(T_{v,e}(M)c)(i,\eta_{\widehat e})=\sum_j M_{ij}c(j,\eta_{\widehat e})$.
\end{definition}

\begin{theorem}[One-edge matrix action on a basis configuration]%
    \label{thm:peps_localIncidentMatrixOp_single}
    \lean{TNLean.PEPS.localIncidentMatrixOp_single}
    \leanok
    \uses{def:peps_localIncidentMatrixOp, def:peps_localVirtualConfigSplitAt}
    For the basis vector $e_{(j,\lambda)}\in\C^{\mathcal V_v}$,
    $T_{v,e}(M)e_{(j,\lambda)}=\sum_i M_{ij}e_{(i,\lambda)}$.
\end{theorem}
\begin{proof}\leanok
    Evaluate Definition~\ref{def:peps_localIncidentMatrixOp} on $(i,\eta_{\widehat e})$
    and split by whether $\eta_{\widehat e}=\lambda$.
\end{proof}

\begin{theorem}[Local tensor form of a one-edge basis action]%
    \label{thm:peps_localTensorMap_localIncidentMatrixOp_single}
    \lean{TNLean.PEPS.localTensorMap_localIncidentMatrixOp_single}
    \leanok
    \uses{thm:peps_localIncidentMatrixOp_single, def:peps_localTensorMap}
    For every $(j,\lambda)$ and physical index $\tau$,
    $\Phi_v(T_{v,e}(M)e_{(j,\lambda)})(\tau)=\sum_i M_{ij}A_v((i,\lambda),\tau)$.
\end{theorem}
\begin{proof}\leanok
    Apply $\Phi_v$ to Theorem~\ref{thm:peps_localIncidentMatrixOp_single} by
    linearity.
\end{proof}

\begin{theorem}[Physical realization of one-edge virtual matrix action]%
    \label{thm:peps_localIncidentMatrixOp_physicalRealization}
    \lean{TNLean.PEPS.localIncidentMatrixOp_physicalRealizationAt}
    \leanok
    \uses{def:peps_localIncidentMatrixOp, def:peps_physRealizeLocalOp, thm:peps_physRealizeLocalOp_spec}
    If $\Phi_v$ is injective, then $O_{T_{v,e}(M)}\in\End(\C^d)$ satisfies
    $O_{T_{v,e}(M)}(\Phi_v(c))=\Phi_v(T_{v,e}(M)c)$ for all
    $c\in\C^{\mathcal V_v}$.
\end{theorem}
\begin{proof}\leanok
    Apply Theorem~\ref{thm:peps_physRealizeLocalOp_spec} to $T=T_{v,e}(M)$.
\end{proof}

\begin{definition}[Split at one incident edge]%
    \label{def:peps_localVirtualConfigSplitAt}
    \lean{TNLean.PEPS.localVirtualConfigSplitAt}
    \leanok
    \uses{def:peps_tensor, def:peps_incident_edge}
    For an incident edge $e\ni v$, the bijection
    $\mathcal V_v\simeq\{0,\ldots,D_e-1\}\times\prod_{f\ni v,\,f\ne e}\{0,\ldots,D_f-1\}$
    writes $\eta\in\mathcal V_v$ as $(\eta(e),\eta_{\widehat e})$.
\end{definition}

\subsection{Edge-centred decomposition of a contraction}

Fixing an edge $e=(u,v)$ splits $V$ into $\{u\}$, $M_e=V\setminus\{u,v\}$, and
$\{v\}$, factoring $\Coeff_A(\sigma)$ for the edge-middle reduction of
Section~\ref{sec:peps_edge_middle}.

\begin{definition}[Middle region of an edge]%
    \label{def:peps_edgeMiddleVertices}
    \lean{TNLean.PEPS.edgeMiddleVertices}
    \leanok
    \uses{def:peps_edge}
    For an edge $e=(u,v)$, the middle region is $M_e=V\setminus\{u,v\}$.
\end{definition}

\begin{theorem}[Edge-centered vertex partition]%
    \label{thm:peps_edgeVertices_union}
    \lean{TNLean.PEPS.edgeVertices_union}
    \leanok
    \uses{def:peps_edgeMiddleVertices}
    For an edge $e=(u,v)$, $V=\{u\}\cup(V\setminus\{u,v\})\cup\{v\}$.
\end{theorem}
\begin{proof}\leanok
    Every $x\in V$ satisfies $x=u$, $x=v$, or $x\in V\setminus\{u,v\}$.
\end{proof}

\begin{theorem}[Nonempty middle region from a cardinal bound]%
    \label{thm:peps_edgeMiddleVertices_nonempty_of_two_lt_card}
    \lean{TNLean.PEPS.edgeMiddleVertices_nonempty_of_two_lt_card}
    \leanok
    \uses{def:peps_edgeMiddleVertices}
    If $2<|V|$, then $V\setminus\{u,v\}\ne\varnothing$ for every edge $e=(u,v)$.
\end{theorem}
\begin{proof}\leanok
    Since $|V\setminus\{u,v\}|=|V|-2$, $2<|V|$ gives $|V\setminus\{u,v\}|>0$.
\end{proof}

\begin{theorem}[Product splitting at an edge]%
    \label{thm:peps_prod_univ_splitAtEdge}
    \lean{TNLean.PEPS.prod_univ_splitAtEdge}
    \leanok
    \uses{def:peps_edgeMiddleVertices, thm:peps_edgeVertices_union}
    For any edge $e=(u,v)$ and $f:V\to M$ into a commutative monoid $M$,
    \begin{align}
        \prod_{x\in V}f(x)
         & =f(u)\left(\prod_{x\in V\setminus\{u,v\}}f(x)\right)f(v).
        \notag
    \end{align}
\end{theorem}
\begin{proof}\leanok
    Split the product over the partition of Theorem~\ref{thm:peps_edgeVertices_union}.
\end{proof}

\begin{theorem}[PEPS coefficient splitting at an edge]%
    \label{thm:peps_stateCoeff_splitAtEdge}
    \lean{TNLean.PEPS.stateCoeff_splitAtEdge}
    \leanok
    \uses{def:peps_stateCoeff, def:peps_edgeMiddleVertices, thm:peps_prod_univ_splitAtEdge}
    For any edge $e=(u,v)$,
    $\Coeff_A(\sigma)=\sum_{\eta}A_u(\eta|_u,\sigma_u)
        (\prod_{x\in V\setminus\{u,v\}}A_x(\eta|_x,\sigma_x))A_v(\eta|_v,\sigma_v)$.
\end{theorem}
\begin{proof}\leanok
    Apply Theorem~\ref{thm:peps_prod_univ_splitAtEdge} to each summand in
    Definition~\ref{def:peps_stateCoeff}.
\end{proof}

\begin{definition}[Edge-centered boundary data]%
    \label{def:peps_edgeBoundaryConfig}
    \lean{TNLean.PEPS.EdgeBoundaryConfig}
    \leanok
    \uses{def:peps_tensor, def:peps_localVirtualConfigSplitAt}
    For an edge $e=(u,v)$, an edge-boundary configuration is a triple
    $\beta=(a,\lambda,\rho)$ consisting of $a\in\{0,\ldots,D_e-1\}$ on $e$ and
    residual configurations $\lambda$ on $\{f\ni u:f\ne e\}$ and $\rho$ on
    $\{f\ni v:f\ne e\}$.
\end{definition}

\begin{definition}[Boundary data read from a virtual configuration]%
    \label{def:peps_edgeBoundaryOfVirtualConfig}
    \lean{TNLean.PEPS.edgeBoundaryOfVirtualConfig}
    \leanok
    \uses{def:peps_edgeBoundaryConfig, def:peps_virtual_config}
    A global virtual configuration $\eta$ determines the boundary configuration
    $\partial_e(\eta)=(\eta(e),\eta|_{u,\widehat e},\eta|_{v,\widehat e})$.
\end{definition}

\begin{definition}[Boundary compatibility at an edge]%
    \label{def:peps_edgeBoundaryMatches}
    \lean{TNLean.PEPS.edgeBoundaryMatches}
    \leanok
    \uses{def:peps_edgeBoundaryConfig, def:peps_virtual_config}
    A global virtual configuration $\eta$ is compatible with
    $\beta=(a,\lambda,\rho)$ if $\partial_e(\eta)=\beta$, i.e., $\eta(e)=a$,
    $\eta|_{u,\widehat e}=\lambda$, and $\eta|_{v,\widehat e}=\rho$.
\end{definition}

\begin{definition}[Middle configurations over fixed boundary data]%
    \label{def:peps_edgeMiddleConfig}
    \lean{TNLean.PEPS.EdgeMiddleConfig}
    \leanok
    \uses{def:peps_edgeBoundaryMatches}
    For fixed boundary data $\beta$, the middle configuration space is the
    fiber $\{\eta:\partial_e(\eta)=\beta\}$.
\end{definition}

\begin{definition}[Virtual configurations fibred by edge boundary data]%
    \label{def:peps_virtualConfigEquivEdgeBoundary}
    \lean{TNLean.PEPS.virtualConfigEquivEdgeBoundary}
    \leanok
    \uses{def:peps_edgeBoundaryOfVirtualConfig, def:peps_edgeMiddleConfig}
    Fiberwise decomposition by $\partial_e$ identifies global virtual
    configurations $\eta$ with pairs $(\beta,\eta_{M_e})$ where
    $\partial_e(\eta_{M_e})=\beta$.
\end{definition}

\begin{definition}[Left endpoint local configuration from boundary data]%
    \label{def:peps_edgeLeftLocalConfig}
    \lean{TNLean.PEPS.edgeLeftLocalConfig}
    \leanok
    \uses{def:peps_edgeBoundaryConfig, def:peps_localVirtualConfigSplitAt}
    From $\beta=(a,\lambda,\rho)$, the left endpoint local configuration
    $\beta_u\in\mathcal V_u$ is reconstructed from $(a,\lambda)$ via
    Definition~\ref{def:peps_localVirtualConfigSplitAt}.
\end{definition}

\begin{definition}[Right endpoint local configuration from boundary data]%
    \label{def:peps_edgeRightLocalConfig}
    \lean{TNLean.PEPS.edgeRightLocalConfig}
    \leanok
    \uses{def:peps_edgeBoundaryConfig, def:peps_localVirtualConfigSplitAt}
    From $\beta=(a,\lambda,\rho)$, the right endpoint local configuration
    $\beta_v\in\mathcal V_v$ is reconstructed from $(a,\rho)$ via
    Definition~\ref{def:peps_localVirtualConfigSplitAt}.
\end{definition}

\begin{theorem}[Left endpoint reconstruction for compatible boundary data]%
    \label{thm:peps_edgeLeftLocalConfig_eq_of_boundaryMatches}
    \lean{TNLean.PEPS.edgeLeftLocalConfig_eq_of_boundaryMatches}
    \leanok
    \uses{def:peps_edgeBoundaryMatches, def:peps_edgeLeftLocalConfig}
    If $\partial_e(\eta)=\beta$, then $\eta|_u=\beta_u$.
\end{theorem}

\begin{theorem}[Right endpoint reconstruction for compatible boundary data]%
    \label{thm:peps_edgeRightLocalConfig_eq_of_boundaryMatches}
    \lean{TNLean.PEPS.edgeRightLocalConfig_eq_of_boundaryMatches}
    \leanok
    \uses{def:peps_edgeBoundaryMatches, def:peps_edgeRightLocalConfig}
    If $\partial_e(\eta)=\beta$, then $\eta|_v=\beta_v$.
\end{theorem}

\begin{definition}[Boundary data with an inserted edge matrix]%
    \label{def:peps_edgeInsertedBoundaryConfig}
    \lean{TNLean.PEPS.EdgeInsertedBoundaryConfig}
    \leanok
    \uses{def:peps_tensor, def:peps_localVirtualConfigSplitAt}
    For an edge $e=(u,v)$, inserted-edge boundary data are quadruples
    $\beta=(a,b,\lambda,\rho)$ with $a,b\in\{0,\ldots,D_e-1\}$ on $e$ and
    residual configurations $\lambda$ at $u$ and $\rho$ at $v$.
\end{definition}

\begin{definition}[Diagonal inserted-edge boundary data]%
    \label{def:peps_edgeBoundaryToInsertedBoundaryConfig}
    \lean{TNLean.PEPS.edgeBoundaryToInsertedBoundaryConfig}
    \leanok
    \uses{def:peps_edgeBoundaryConfig, def:peps_edgeInsertedBoundaryConfig}
    Ordinary boundary data $\beta=(a,\lambda,\rho)$ map to diagonal
    inserted-edge boundary data $\Delta(\beta)=(a,a,\lambda,\rho)$.
\end{definition}

\begin{definition}[Right-endpoint inserted-boundary reindexing]%
    \label{def:peps_edgeBoundaryRightIndexEquivInsertedBoundaryConfig}
    \lean{TNLean.PEPS.edgeBoundaryRightIndexEquivInsertedBoundaryConfig}
    \leanok
    \uses{def:peps_edgeBoundaryConfig, def:peps_edgeInsertedBoundaryConfig}
    Pairs $((a,\lambda,\rho),b)$ of ordinary boundary data and a right bond
    index $b$ identify bijectively with inserted-edge boundary data
    $(a,b,\lambda,\rho)$.
\end{definition}

\begin{definition}[Left-endpoint inserted-boundary reindexing]%
    \label{def:peps_edgeBoundaryLeftIndexEquivInsertedBoundaryConfig}
    \lean{TNLean.PEPS.edgeBoundaryLeftIndexEquivInsertedBoundaryConfig}
    \leanok
    \uses{def:peps_edgeBoundaryConfig, def:peps_edgeInsertedBoundaryConfig}
    Pairs $((b,\lambda,\rho),a)$ of ordinary boundary data and a left bond
    index $a$ identify bijectively with inserted-edge boundary data
    $(a,b,\lambda,\rho)$.
\end{definition}

\begin{definition}[Diagonal inserted-boundary configurations]%
    \label{def:peps_edgeDiagonalInsertedBoundaryConfig}
    \lean{TNLean.PEPS.EdgeDiagonalInsertedBoundaryConfig}
    \leanok
    \uses{def:peps_edgeInsertedBoundaryConfig}
    The diagonal inserted-boundary configurations are the subtype
    $\{(a,b,\lambda,\rho):a=b\}$.
\end{definition}

\begin{definition}[Diagonal boundary reindexing]%
    \label{def:peps_edgeBoundaryConfigEquivDiagonalInsertedBoundaryConfig}
    \lean{TNLean.PEPS.edgeBoundaryConfigEquivDiagonalInsertedBoundaryConfig}
    \leanok
    \uses{def:peps_edgeBoundaryConfig, def:peps_edgeDiagonalInsertedBoundaryConfig, def:peps_edgeBoundaryToInsertedBoundaryConfig}
    The diagonal map $\Delta:(a,\lambda,\rho)\mapsto(a,a,\lambda,\rho)$ is a
    bijection onto the diagonal inserted-boundary configurations.
\end{definition}
